\chapter{Future Work}
\label{sec:future-work}

\paragraph{Human evaluation} is the most urgent item, because only human judgements of the simplified outputs can resolve the document-level metric disagreement.
\Gls{dsari} ranks the fine-tuned document model above the prompted 7\gls{billion}, and \gls{lens} ranks them the other way round (\cref{sec:document-results}).
\Gls{lens} was trained on sentence-level human ratings \parencite[16385--16386]{maddela-etal-2023-lens}, so on documents it is applied out of domain and cannot settle the disagreement it creates (\cref{sec:limitations}).

A study of this disagreement needs two systems, three raters, and 30 to 50 documents.
Two systems suffice because the disagreement concerns only these two, and with three raters no single rater decides an item.
The number of documents is limited by reading load: at a corpus mean of about 180 tokens per source document (\cref{sec:granularity-findings}), 30 to 50 documents already give each rater 60 to 100 outputs to read against their sources.
The raters would score meaning preservation, fluency, and simplicity, the three standard dimensions \parencite[147]{alva-manchego-etal-2020-survey}, and the study would report inter-annotator agreement.

The samples already exist: the paired document sets are on disk with a verified fingerprint, and the \num{130} labelled audit items carry content-hash identifiers.
The four error cases of \cref{sec:qualitative} (reciting the corpus over the input, over-deletion, fluent invention on technical prose, declining to edit) give the raters a taxonomy to rate against.
What is missing are the participants, an ethics and consent process, and rater training.
The ratings could later serve as training data for a learned metric at document level, which \gls{lens} is not.

\paragraph{A polarity check in the change guard} would address the five negation deletions that inverted a sentence's meaning (\cref{sec:deployment-results}).
The obvious implementation is a negation stop-list, which was not built: like the two-entry deny-list of \cref{sec:limitations}, it would ship without a measured false-positive rate.
The measurement that matters first is how often such inversions occur.
Run on more real prose, the detection method of \cref{app:inversions} gives a lower bound per thousand sentences and needs no \gls{gpu}.

\pagebreak

\paragraph{A language gate} should come before any multilingual training.
Both checkpoints are English-only, and nothing checks the language of the text sent to them.
On one German page, the served outputs contained two deleted negations and lowercased non-words, and no guard fired.
A gate that detects a unit's language and declines non-English text would make the system's behaviour match its documented English-only scope; it would not add multilingual capability.

Multilingual training comes second because it raises a question the English pipeline never had to answer: which simplified register to target.
German, for example, codifies two, \emph{Leichte Sprache} (``Easy Language'') and \emph{Einfache Sprache} (``Plain Language''), each with its own rule set and its own audience \parencite[21, 42]{stodden2026automatic}.

\paragraph{Natural-case document training} would make the normalisation layer unnecessary.
The document model consumes and produces the corpus's lowercased, pre-tokenised convention, so the deployed system normalises its input and de-normalises its output (\cref{sec:unit-of-simplification}).
This pair is a handwritten approximation of the corpus's tokenisation, and it caused two defects: it shifted a section's paragraphs out of step (\cref{sec:granularity-findings}) and produced the German non-words of \cref{sec:deployment-results}.
Training on natural-case, normally punctuated text would remove the layer and the failure modes that come with it.

A cheaper first step is to measure what the layer contributes.
The input-convention comparison of \cref{tab:conventions} rests on one hand-constructed document; repeating it on many natural-case documents, with and without normalisation, needs new generations but no retraining.

\paragraph{Cohesion} was not measured, although over-deletion is the deployed checkpoint's dominant quality failure on long documents (\cref{sec:qualitative}).
Deletion can separate pronouns from their antecedents or connectives from the clauses they link, raising the question of whether cohesion decreases as compression increases.
Future work could measure cohesion using \gls{taaco}-style indices \parencite[1227]{crossley-etal-2016-taaco} across source, reference, and system outputs.
Since global cohesion indices could predict coherence judgements in the validation of \gls{taaco} \parencite[1227, 1233--1234]{crossley-etal-2016-taaco}, global and local measures should both be reported, although the global indices compare paragraphs \parencite[1227]{crossley-etal-2016-taaco}, which the D-Wikipedia documents do not mark (\cref{sec:granularity-findings}).

If cohesion decreases with compression, it could become a target for controllable simplification. Conditioning on cohesion with a control token, following \gls{access} \parencite[4691]{martin-etal-2020-controllable}, would require retraining and a validated cohesion metric.
SimDoc's joint objective \parencite[1]{vasquezrodriguez2024simple} provides a related existing approach.

\paragraph{Two filter changes} were quantified but not made, because each also has a cost.
Dropping the semicolon from the code heuristic would send all \num{102} items it matched, not only the 17 that were wanted (\cref{sec:rq2}); a narrower rule would have to separate those 17 from the other 85.

Raising the minimum word count from four to six would stop \num{647} of the \num{2222} sentence-cut items from being sent, more than any other change measured.
It would also drop \num{349} rewrites of short text whose value was never labelled (\cref{sec:deployment-results}), so labelling them would decide whether the change is worth making.

\paragraph{In-browser inference} would remove the local backend constraints, and runtime measurements show which models it suits.
The fine-tuned sentence and document checkpoints take \(\approx\)\qty{20}{\second} per page, while the prompted 3\gls{billion} can take up to \qty{2}{\minute} (\cref{sec:runtime}).
Therefore, should investigate whether \gls{bart}-based \gls{seqseq} models and prompted \glspl{llm}, if they could run in the browser at a usable speed, and how their runtime compares with the local backend.

\paragraph{Audience control} exists here only as an instruction to the prompted models, and whether it serves the audiences it names was not tested.
Four of its five audiences have no reference data (\cref{sec:limitations}), and their demonstrations are hand-authored (\cref{app:demos}).

A human study with readers from these audiences, could show whether the audience parameter helps them specifically.
With the help of generative AI models, a corpus of such audience-conditioned examples could be built.
A separate model could also be fine-tuned on them, so that trained and instructed audience control could be compared directly.
Fine-tuning the prompted model itself, for example parameter-efficiently, is excluded, because it enters \cref{rq:comparison} as an instructed rather than a trained system, and fine-tuning it would remove that distinction.
